\documentclass[12pt]{article}
\usepackage[utf8]{inputenc}
\usepackage[T1]{fontenc}
\usepackage[margin=1in]{geometry}
\usepackage{setspace}
\usepackage{xcolor}
\usepackage[colorlinks=true, urlcolor=blue!60!black, linkcolor=blue!60!black]{hyperref}
\setlength{\parindent}{0em}
\setlength{\parskip}{0.45em}
\renewcommand{\baselinestretch}{1.0}
\begin{document}
\begin{center}
{\Large \textit{Sveen v. Melin}, 584 U.S. 811 (2018)\par}
\vspace{0.3em}
{\normalsize Case Brief --- One-Page Edition\par}
\end{center}

\textbf{Memory Jogger.} Revocation-on-divorce and the Contracts Clause (\href{https://supreme.justia.com/cases/federal/us/584/16-1432/}{584 U.S.\ 811}).

\textbf{Facts.} Mark Sveen named his wife, Kaye Melin, beneficiary of his life insurance in 1998. In 2002 Minnesota extended its revocation-on-divorce rule to life insurance. The couple divorced in 2007, Mark never changed the designation, and after his death Melin and his children both claimed the proceeds.

\textbf{Procedural History.} The district court ruled for the children; the Eighth Circuit reversed under the Contracts Clause.

\textbf{Votes.} Reversed, 8--1 (Kagan for the Court; Gorsuch dissenting).

\textbf{Holding.} Retroactive application of the statute does not violate the Contracts Clause because it does not substantially impair the contract, so means-ends review is unnecessary.

\textbf{Reasoning.} The law reflects likely intent, since ``the average Joe does not want his ex inheriting what he leaves behind''; it does no more than a divorce court could do; and the policyholder can undo it ``with the stroke of a pen,'' like filing laws upheld since the 1800s. Gorsuch argued the Clause bars any retroactive impairment and that naming the beneficiary is the ``whole point'' of the contract.

\textbf{Examples.} \emph{Same-Side:} payable-on-death accounts. \emph{Opposite-Side:} retroactively voiding non-competes. \emph{Fence-Sitter:} a revocation rule undone only by court petition.

\textbf{Critique.} The rationale assumes policyholders know the law changed, and step one absorbs tailoring.

\textbf{Key Quotation.} ``In answering that question, the Court has considered the extent to which the law undermines the contractual bargain, interferes with a party's reasonable expectations, and prevents the party from safeguarding or reinstating his rights.'' --- Justice Kagan, 584 U.S.\ at 819.

\end{document}
